\section*{Project Assignment}
\label{sec:project_assignment}

\begin{table}[htbp]
    \centering
    \renewcommand{\arraystretch}{1.2}
    \begin{tabular}{@{}ll@{}}
        \toprule
        \textbf{Candidate} & Anders Haarberg Eriksen \\
        \addlinespace
        \textbf{Programme} & Cybernetics and Robotics \\
        \addlinespace
        \textbf{Project title} &
        \begin{tabular}[t]{@{}l@{}}
            Adaptive Control and Parameter Estimation for Motion Control \\
            in a Tendon-Driven Hand Exoskeleton
        \end{tabular} \\
        \addlinespace
        \textbf{Norwegian title} &
        \begin{tabular}[t]{@{}l@{}}
            Adaptiv regulering og parameterestimering for bevegelsesstyring \\
            i et senedrevet håndeksoskjelett
        \end{tabular} \\
        \addlinespace
        \textbf{Supervisors} & Øyvind Stavdahl, Devon Rolhf \\
        \addlinespace
        \textbf{Institution} &
        \begin{tabular}[t]{@{}l@{}}
            Department of Engineering Cybernetics, NTNU, \\
            in collaboration with the N-squared Lab, FAU, Erlangen, Germany
        \end{tabular} \\
        \bottomrule
    \end{tabular}
\end{table}

Modern hand exoskeletons require precise and reliable control of the relevant mechanical states to enable natural, predictable, and safe interaction with the environment. However, the current implementation of the present glove--cable system introduces nonlinearities, uncertain mechanical parameters, and inconsistent finger positioning, making accurate control difficult. To address these challenges, an adaptive control framework is desired that can estimate system parameters in real time, compensate for nonlinearities, and ensure stable behavior during both free motion and object interaction. 

The goal of this project is to design an adaptive controller for a \mbox{6-DOF} hand exoskeleton, capable of regulating both torque and fingertip positioning. This includes identifying relevant mechanical parameters and creating a nonlinear model of the system. Additionally, modifications to the glove mechanics and integration of fingertip sensors could enable better positioning and torque estimation and facilitate switching between control modes and should therefore be investigated. The final system should allow reliable, intuitive control of finger motion while providing insight into system properties such as friction and internal forces.

The work should thus include the following:
\begin{itemize}
	\item Development of a nonlinear model of a single finger to support adaptive control.
    \item Design of an adaptive control architecture for both torque output and finger positioning.
    \item Estimation of mechanical parameters and internal forces.
    \item Evaluation of fingertip sensors to support mode switching and refined control.
    \item Documentation and of the system architecture, experiments, and control performance.
\end{itemize}